% GENERATED FILE. Edit canonical lessons, then run npm run generate:books.
% canonical-source-hash: fnv1a64:335da39f17d00b97
% canonical-lessons: PT-C501-jardinagem, PT-C501-semear, PT-C501-regar, PT-C501-regador, PT-C501-vaso

\chapter{Mãos na terra}
\label{ch:pt-maos-na-terra}
\hlchaptermodality{501}

\begin{chapteropening}
\textbf{By the end of this chapter:} \emph{I can talk about gardening.}

Complete the last lesson of chapter 501: I can talk about gardening.
\end{chapteropening}

\section[a jardinagem]{a jardinagem — gardening}

\label{lesson:PT-C501-jardinagem}



\begin{quote}
Before the new one: say the Portuguese for to knit, then the Portuguese for to embroider.
\end{quote}

\subsection*{You'll want to know: a jardinagem}
\textbf{a jardinagem} — \textquotedblleft{}gardening\textquotedblright{}.

\begin{quote}
A jardinagem é o passatempo preferido do meu pai.
\end{quote}

\begin{quote}
\emph{Gardening is my father's favourite hobby.}
\end{quote}

Built on \emph{o jardim}.

\begin{cousinweb}
\emph{jardinagem} comes from French \emph{jardinage} \textquotedblleft{}gardening\textquotedblright{} (EtymDB), built on \emph{jardin} \textquotedblleft{}garden\textquotedblright{}, a word of Germanic origin.
\end{cousinweb}

\begin{grammarlens}[title={What you've built}]
One more word for telling, explaining and describing at B1.
\end{grammarlens}

\subsection*{Your turn}
\begin{itemize}
\raggedright
  \item \emph{Say it:} \emph{a jardinagem}
  \item \emph{Say it:} \emph{a jardinagem}, once more
  \item \emph{Say it:} a sentence of your own with \emph{a jardinagem}
  \item \emph{Recall:} say the Portuguese for to knit, then the Portuguese for to embroider, then say \emph{a jardinagem} again
\end{itemize}

\begin{tcolorbox}[breakable,colback=teal!4,colframe=teal!35!black,title={Before you move on}]
What does a jardinagem mean? (\textquotedblleft{}Gardening\textquotedblright{}.) Say it once more.
\end{tcolorbox}
\section[semear]{semear — to sow}

\label{lesson:PT-C501-semear}



\begin{quote}
Before the new one: say the Portuguese for to embroider, then the Portuguese for gardening.
\end{quote}

\subsection*{You'll want to know: semear}
\textbf{semear} — \textquotedblleft{}to sow\textquotedblright{}.

\begin{quote}
Em março semeamos tomates.
\end{quote}

\begin{quote}
\emph{In March we sow tomatoes.}
\end{quote}

Present: \emph{semeio}, \emph{semeias}, \emph{semeia}.

\begin{cousinweb}
\emph{semear} is inherited, through Old Portuguese (EtymDB), from Latin \emph{sēmināre} \textquotedblleft{}to sow\textquotedblright{}, built on \emph{sēmen} \textquotedblleft{}seed\textquotedblright{} (Lewis \& Short).
\end{cousinweb}

\begin{grammarlens}[title={What you've built}]
One more word for telling, explaining and describing at B1.
\end{grammarlens}

\subsection*{Your turn}
\begin{itemize}
\raggedright
  \item \emph{Say it:} \emph{semear}
  \item \emph{Say it:} \emph{semear}, once more
  \item \emph{Say it:} a sentence of your own with \emph{semear}
  \item \emph{Recall:} say the Portuguese for to embroider, then the Portuguese for gardening, then say \emph{semear} again
\end{itemize}

\begin{tcolorbox}[breakable,colback=teal!4,colframe=teal!35!black,title={Before you move on}]
What does semear mean? (\textquotedblleft{}To sow\textquotedblright{}.) Say it once more.
\end{tcolorbox}
\section[regar]{regar — to water}

\label{lesson:PT-C501-regar}



\begin{quote}
Before the new one: say the Portuguese for gardening, then the Portuguese for to sow.
\end{quote}

\subsection*{You'll want to know: regar}
\textbf{regar} — \textquotedblleft{}to water\textquotedblright{}.

\begin{quote}
Rego as plantas todas as manhãs.
\end{quote}

\begin{quote}
\emph{I water the plants every morning.}
\end{quote}

\emph{Regar o jardim} — to water the garden.

\begin{cousinweb}
\emph{regar} is inherited from Latin \emph{rigāre} \textquotedblleft{}to water, to wet\textquotedblright{} (EtymDB).
\end{cousinweb}

\begin{grammarlens}[title={What you've built}]
One more word for telling, explaining and describing at B1.
\end{grammarlens}

\subsection*{Your turn}
\begin{itemize}
\raggedright
  \item \emph{Say it:} \emph{regar}
  \item \emph{Say it:} \emph{regar}, once more
  \item \emph{Say it:} a sentence of your own with \emph{regar}
  \item \emph{Recall:} say the Portuguese for gardening, then the Portuguese for to sow, then say \emph{regar} again
\end{itemize}

\begin{tcolorbox}[breakable,colback=teal!4,colframe=teal!35!black,title={Before you move on}]
What does regar mean? (\textquotedblleft{}To water\textquotedblright{}.) Say it once more.
\end{tcolorbox}
\section[o regador]{o regador — watering can}

\label{lesson:PT-C501-regador}



\begin{quote}
Before the new one: say the Portuguese for to sow, then the Portuguese for to water.
\end{quote}

\subsection*{You'll want to know: o regador}
\textbf{o regador} — \textquotedblleft{}watering can\textquotedblright{}.

\begin{quote}
Encho o regador na torneira do quintal.
\end{quote}

\begin{quote}
\emph{I fill the watering can at the tap in the yard.}
\end{quote}

Built on \emph{regar} with \emph{-dor}, the ending for a tool.

\begin{cousinweb}
\emph{regador} is built on \emph{regar} \textquotedblleft{}to water\textquotedblright{}, inherited from Latin \emph{rigāre} (EtymDB).
\end{cousinweb}

\begin{grammarlens}[title={What you've built}]
One more word for telling, explaining and describing at B1.
\end{grammarlens}

\subsection*{Your turn}
\begin{itemize}
\raggedright
  \item \emph{Say it:} \emph{o regador}
  \item \emph{Say it:} \emph{o regador}, once more
  \item \emph{Say it:} a sentence of your own with \emph{o regador}
  \item \emph{Recall:} say the Portuguese for to sow, then the Portuguese for to water, then say \emph{o regador} again
\end{itemize}

\begin{tcolorbox}[breakable,colback=teal!4,colframe=teal!35!black,title={Before you move on}]
What does o regador mean? (\textquotedblleft{}Watering can\textquotedblright{}.) Say it once more.
\end{tcolorbox}
\section[o vaso]{o vaso — flowerpot; vase}

\label{lesson:PT-C501-vaso}



\begin{quote}
Before the new one: say the Portuguese for to water, then the Portuguese for watering can.
\end{quote}

\subsection*{You'll want to know: o vaso}
\textbf{o vaso} — \textquotedblleft{}flowerpot; vase\textquotedblright{}.

\begin{quote}
Tenho um vaso de flores na varanda.
\end{quote}

\begin{quote}
\emph{I have a pot of flowers on the balcony.}
\end{quote}

In Portugal mainly a pot for plants; cut flowers go in \emph{uma jarra}.

\begin{cousinweb}
\emph{vaso} is inherited, through Old Portuguese, from Latin \emph{vāsum} \textquotedblleft{}a vessel\textquotedblright{} (EtymDB).
\end{cousinweb}

\begin{grammarlens}[title={What you've built}]
One more word for telling, explaining and describing at B1.
\end{grammarlens}

\subsection*{Your turn}
\begin{itemize}
\raggedright
  \item \emph{Say it:} \emph{o vaso}
  \item \emph{Say it:} \emph{o vaso}, once more
  \item \emph{Say it:} a sentence of your own with \emph{o vaso}
  \item \emph{Recall:} say the Portuguese for to water, then the Portuguese for watering can, then say \emph{o vaso} again
\end{itemize}

\begin{tcolorbox}[breakable,colback=teal!4,colframe=teal!35!black,title={Before you move on}]
What does o vaso mean? (\textquotedblleft{}Flowerpot; vase\textquotedblright{}.) Say it once more.
\end{tcolorbox}
